% !TeX root = ../../Book1.tex
\section{Frostman 引理}
\label{b1:sec:1504}

在 Cantor 集上，基本区间已经告诉我们应把质量放在哪里。
一般紧集没有现成的分支结构，却仍可以提出同一要求：半径为 \(r\) 的球至多承载常数倍的 \(r^s\) 质量。
质量分布原理说明，这种要求足以给出维数下界。
\term[frostman-yinli]{Frostman 引理}[Frostman's lemma]则回答反方向的问题：集合的 Hausdorff 测度为正时，是否总能找到这样的质量分布？
本节对欧氏空间中的紧集给出肯定回答，并在本节内证明构造所需的紧集概率测度序列紧性。

% 来源范围：Fremlin2016Measure 471Xd(ii)只作为Frostman型结论的对照，
% 不把该习题的附加假设误当成一般紧集定理的完整证明。
% 下文有限树递推、割集下界、原子测度弱极限及能量估计均在本书内完整展开。
% Mattila1995Geometry与Falconer2014Fractal在本节只作专题伴读推荐。

\subsection{质量分布原理回顾与统一增长条件}
\label{b1:subsec:150401}

\cref{b1:lem:mass-distribution}已经正式给出了质量分布原理及其证明。
本小节只回顾其中的覆盖机制，并强调 Frostman 理论所需估计必须在中心、常数与尺度上统一。
设 \(s>0\)，\(\mu\) 是集中在 Borel 集 \(E\subset\mathbb R^d\) 上的有限正测度，
\(\mu(E)>0\)。若所有足够小的球都满足
\[
 \mu\bigl(B(x,r)\bigr)\leq A r^s,
\]
则任意细覆盖都必须付出一定的直径幂和。
事实上，对与 \(E\) 相交且直径为 \(a>0\) 的覆盖集 \(U\)，取 \(x\in U\cap E\)，
就有 \(U\subset\overline B(x,a)\)，从而其质量至多为 \(Aa^s\)。
这里不要求 \(U\) 可测：应写作 \(\mu^*(U\cap E)\)，再用包住它的可测球估计。
这正是\cref{b1:lem:mass-distribution}的覆盖论证。

值得注意的是不等式的方向。若希望证明集合至少有 \(s\) 维，
就应限制小球不能装入过多的质量，而不是要求小球质量至少为 \(r^s\)。
一旦总质量固定，单个小集合能承载的质量越少，覆盖整个集合所需的集合就越多。
常数 \(A\) 只改变最后的测度下界，不改变维数。

若估计只在 \(\mu\)-几乎处处中心成立，也常可先限制到一个正质量子集上。
但必须检查两个统一性问题：常数是否统一，以及允许的半径范围是否统一。
例如，若对每个 \(x\) 都有某个 \(A_x<\infty\) 和 \(r_x>0\)，
可按 \(A_x\leq n\)、\(r_x\geq1/n\) 分组；在分组可测的条件下，
总有一组仍有正质量。仅写“对每个点有局部估计”而不作这一步，不能直接套用质量分布原理。

\subsection{Frostman 测度}
\label{b1:subsec:150402}

\begin{definition}\label{b1:def:frostman-measure}
设 \(s\geq0\)，\(K\subset\mathbb R^d\) 为非空紧集。
若非零有限正 Borel 测度 \(\mu\) 集中在 \(K\) 上，并存在 \(A<\infty\)，使
\[
 \mu\bigl(B(x,r)\bigr)\leq A r^s
 \quad(x\in\mathbb R^d,\ r>0),
\]
则称 \(\mu\) 为 \(K\) 上的一个 \term[frostman-cedu]{Frostman 测度}[Frostman measure]，
其中 \(s\) 是所用的增长指数。
\end{definition}

“集中在 \(K\) 上”只要求 \(\mu(\mathbb R^d\setminus K)=0\)，
不要求 \(\operatorname{supp}\mu=K\)。
由于总质量有限，除以 \(\mu(K)\) 即可改成概率测度，只需相应调整 \(A\)。
当 \(s>0\) 时，这样的测度没有原子，因为令 \(r\downarrow0\) 就得 \(\mu(\{x\})=0\)。
当 \(s=0\) 时，非空紧集上的任意 Dirac 测度均满足要求。

\begin{lemma}\label{b1:lem:hausdorff-content-positive}
若 \(s>0\)，则对任意集合 \(E\subset\mathbb R^d\)，
\[
 \mathcal H^s(E)>0\quad\Longleftrightarrow\quad
 \mathcal H^s_\infty(E)>0.
\]
\end{lemma}
\begin{proof}
这正是\cref{b1:prop:hausdorff-borel-hull}中零测性等价关系的逆否形式：
无尺度限制的代价能够任意小时，各个覆盖集的直径也被迫任意小。
\end{proof}

下面的构造不把 \(\mathcal H^s|_K\) 归一化，因为 \(\mathcal H^s(K)\) 可能为无穷。
我们先在有限层的二进立方体上分配质量，再把层数趋于无穷。

这里的二进立方体是半开立方体 \(Q_0\) 及反复等分各边所得的子立方体；
每个立方体有 \(2^d\) 个互不相交的子立方体。记其边长为 \(\ell(Q)\)。
半开约定使同层立方体真正构成分割，点不会因落在公共边界而被重复计数。

\cref{b1:fig:frostman-dyadic-tree}示意有限树上的两个约束：质量从根向下守恒，
每个结点所得质量又不超过递推给出的容量 \(F_N(Q)\)。末层质量最终放在
\(K\) 中选定的点上，形成离散测度。

\begin{figure}[htbp]
\centering
\begin{tikzpicture}[>=stealth,line cap=round,line join=round,
  box/.style={draw,rounded corners,align=center,inner sep=4pt}]
  \node[box] (root) at (0,2.2) {$Q_0$\\$m(Q_0)=a$};
  \node[box] (left) at (-3.4,0.5) {$Q'$\\$m(Q')\leq F_N(Q')$};
  \node at (0,0.5) {$\cdots$};
  \node[box] (right) at (3.4,0.5) {$Q''$\\$m(Q'')\leq F_N(Q'')$};
  \node[box] (ll) at (-4.5,-1.35) {$R_1$\\$m(R_1)\delta_{x_{R_1}}$};
  \node[box] (lr) at (-2.3,-1.35) {$R_2$\\$m(R_2)\delta_{x_{R_2}}$};
  \node at (2.25,-1.35) {$\cdots$};
  \node[box] (rr) at (4.15,-1.35) {$R$\\$m(R)\delta_{x_R}$};
  \draw[->,thick] (root)--(left);
  \draw[->,thick] (root)--(right);
  \draw[->,thick] (left)--(ll);
  \draw[->,thick] (left)--(lr);
  \draw[->,thick] (right)--(rr);
  \node[align=center] at (0,-1.35) {逐层分配\\质量守恒};
\end{tikzpicture}
\caption[二进质量分配]{二进质量分配。根质量 \(a\) 逐层向子立方体分配；每个结点受容量 \(F_N\) 控制，末层质量则变成位于 \(K\) 上的原子。图中只画出部分分支。}
\label{b1:fig:frostman-dyadic-tree}
\end{figure}

下面把构造中唯一需要的紧性输入单独列出。这样本节不依赖\chapref{b1:ch:12}的完整弱收敛理论；
该章给出的是波兰空间、紧性与 Portmanteau 判据的更一般版本。

\begin{lemma}[紧集上概率测度的序列紧性]\label{b1:lem:compact-probability-subsequence}
设 \(K\subset\mathbb R^d\) 为非空紧集，\((\nu_n)\) 是 \(K\) 上的 Borel 概率测度序列。
则存在子列 \((\nu_{n_j})\) 与 \(K\) 上的 Borel 概率测度 \(\nu\)，使
\[
 \int_K f\,\mathrm d\nu_{n_j}\longrightarrow\int_K f\,\mathrm d\nu
 \quad\bigl(f\in C(K;\mathbb R)\bigr).
\]
此外，对每个在 \(K\) 中相对开的集合 \(G\)，都有
\[
 \nu(G)\leq\liminf_{j\to\infty}\nu_{n_j}(G).
\]
\end{lemma}
\begin{proof}
先说明 \(C(K)\) 可分。对每个 \(n\geq1\)，取有限的 \(1/n\)-网
\(\{q_{n,1},\ldots,q_{n,N_n}\}\)，并令
\[
 w_{n,k}(x)=\max\{0,2/n-|x-q_{n,k}|\},
 \qquad
 \phi_{n,k}(x)=\frac{w_{n,k}(x)}{\sum_iw_{n,i}(x)}.
\]
分母处处为正，且 \(\sum_k\phi_{n,k}=1\)。连续函数在紧集上一致连续，故
\(\sum_k f(q_{n,k})\phi_{n,k}\) 随 \(n\to\infty\) 一致逼近 \(f\)。
再把有限个系数换成有理数，就得到 \(C(K;\mathbb R)\) 的可数稠密族。

由于 \(K\) 紧，\(C(K;\mathbb R)=C_0(K;\mathbb R)\)。由\cref{b1:cor:positive-c0-representation}，
概率测度恰对应于 \(C(K;\mathbb R)^*\) 中满足
\[
 L(1)=1,\qquad L(f)\geq0\quad(f\geq0)
\]
的泛函；它们都位于对偶单位球中。这些条件在弱-*拓扑下闭。
\cref{b1:thm:banach-alaoglu}、\cref{b1:thm:weak-star-metrizable}与 \(C(K)\) 的可分性说明，
这族泛函是紧致可度量空间，因而每个序列都有弱-*收敛子列。
再用同一个表示定理识别极限泛函，就得到概率测度 \(\nu\) 及题设的积分收敛。

最后设 \(G\subset K\) 相对开。若 \(G=K\)，所需不等式由总质量为 \(1\) 得到；
否则令
\[
 f_m(x)=\min\{1,m\,\operatorname{dist}(x,K\setminus G)\}.
\]
则 \(f_m\in C(K)\)、\(0\leq f_m\leq\mathbf1_G\)，且 \(f_m\uparrow\mathbf1_G\)。
于是对每个 \(m\)，
\[
 \int_K f_m\,\mathrm d\nu
 =\lim_{j\to\infty}\int_K f_m\,\mathrm d\nu_{n_j}
 \leq\liminf_{j\to\infty}\nu_{n_j}(G).
\]
令 \(m\to\infty\)，再用单调收敛定理，即得开集不等式。
\end{proof}

\begin{theorem}[Frostman]\label{b1:thm:frostman}
设 \(K\subset\mathbb R^d\) 为非空紧集，\(0<s\leq d\)。
则 \(\mathcal H^s(K)>0\) 当且仅当 \(K\) 上存在增长指数为 \(s\) 的 Frostman 测度。
\end{theorem}
存在性方向有四个转换：先把正覆盖量变成有限二进立方体树上的容量下界；
再从根向下分配质量，得到具有统一立方体控制的离散测度；
随后用前一引理抽取弱收敛子列；最后用开集不等式把球质量控制传给极限。即
\[
 \begin{aligned}
 \mathcal H^s_\infty(K)>0
 &\longrightarrow \text{有限树容量与质量分配}
 \longrightarrow \mu_N(B(x,r))\lesssim r^s\\
 &\longrightarrow \mu_{N_j}\Rightarrow\mu
 \longrightarrow \mu(B(x,r))\lesssim r^s.
 \end{aligned}
\]
\begin{proof}
已有测度时，由 \(\overline B(x,r)\subset B(x,2r)\) 及质量分布原理即得
\(\mathcal H^s(K)>0\)。
以下设 \(\mathcal H^s(K)>0\)，取内部包含 \(K\) 的半开立方体 \(Q_0\)，边长为 \(L\)。
由前一引理，
\[
 a=\frac{\mathcal H^s_\infty(K)}{c_s d^{s/2}}>0.
\]
又因 \(Q_0\) 本身覆盖 \(K\)，有 \(a\leq L^s\)。

固定层数 \(N\)。给第 \(N\) 层立方体规定数值
\[
 F_N(Q)=
 \begin{cases}
  \ell(Q)^s,&Q\cap K\ne\varnothing,\\
  0,&Q\cap K=\varnothing.
 \end{cases}
\]
然后逐层向上定义
\begin{equation}\label{b1:eq:frostman-tree-recursion}
 F_N(Q)=\min\left\{\ell(Q)^s,\ \sum_{Q'\text{ 为 }Q\text{ 的子立方体}}F_N(Q')\right\}.
\end{equation}
这个递推有一个直接的覆盖解释：在 \(Q\) 内，要覆盖所有与 \(K\) 相交的末层立方体，
可以选取 \(Q\) 本身，也可以分别在各个子立方体内继续选取。
对有限层数归纳，\(F_N(Q_0)\) 就是这类选取所能达到的最小边长幂和。
更明确地说，每条通往非空末层立方体的分支至少须选中一个立方体；
若选中了某立方体，就不必再选它的后代。
每一种如此得到的有限族都覆盖 \(K\)，所以
\[
 c_s d^{s/2} F_N(Q_0)\geq\mathcal H^s_\infty(K),
 \quad F_N(Q_0)\geq a.
\]

现在从上向下分配质量。给根立方体质量 \(a\)；
若 \(Q\) 已分得 \(m(Q)\leq F_N(Q)\)，由%
\eqref{b1:eq:frostman-tree-recursion}可把它分成非负数 \(m(Q')\leq F_N(Q')\)，
且所有子立方体所得质量之和恰为 \(m(Q)\)。
这种分配总能完成：依次向子立方体分配不超过其限额的质量，
由于所有限额之和至少为 \(m(Q)\)，不会有剩余质量无处安放。
有限次重复后，在每个与 \(K\) 相交的末层立方体内选一点 \(x_Q\)，令
\[
 \mu_N=\sum_{Q\text{ 在第 }N\text{ 层}}m(Q)\delta_{x_Q},
\]
空立方体不计入和式。于是 \(\mu_N(K)=a\)，
而每个不深于第 \(N\) 层的立方体均满足
\[
 \mu_N(Q)=m(Q)\leq\ell(Q)^s.
\]
这里使用了半开分割，故原子即使位于某些立方体边界上，等式仍然成立。

概率测度 \(\mu_N/a\) 全部集中在紧集 \(K\) 上。
由\cref{b1:lem:compact-probability-subsequence}，可取子列 \((\mu_{N_j}/a)\)
弱收敛到 \(K\) 上的概率测度；把极限乘以 \(a\)，并视为
\(\mathbb R^d\) 上的测度，记为 \(\mu\)。
因此 \(\mu(K)=a>0\)。还须证明小球质量估计在极限中保留；
不能直接对半开立方体使用闭集版本的 Portmanteau 不等式。

若 \(0<r<L\)，选取一层，使其边长 \(\ell\) 满足 \(r\leq\ell<2r\)。
球 \(B(x,r)\) 在每个坐标方向至多遇到四个该层区间，
故它至多与 \(4^d\) 个该层立方体相交。只要 \(N\) 不小于这一层数，就有
\[
 \mu_N\bigl(B(x,r)\bigr)\leq4^d\ell^s\leq4^d2^s r^s.
\]
对 \(K\cap B(x,r)\) 使用\cref{b1:lem:compact-probability-subsequence}中的开集不等式，得
\[
 \mu\bigl(B(x,r)\bigr)
 \leq\liminf_{j\to\infty}\mu_{N_j}\bigl(B(x,r)\bigr)
 \leq4^d2^s r^s,
\]
其中只需取 \(j\) 充分大，使 \(N_j\) 不小于所选层数。
若 \(r\geq L\)，则 \(\mu\bigl(B(x,r)\bigr)\leq a\leq L^s\leq r^s\)。
故 \(\mu\) 满足定义中的统一增长估计。
\end{proof}

紧性在这段证明中有明确的位置：它使各层构造的质量无法在弱极限中逃走。
另一方面，各层的质量分配不必彼此相容；相容性由取弱极限解决。
把这两件事分开，往往比直接在一棵无限树上规定所有分支的质量更方便。
关于更一般集合上的版本及其与势论的关系，可继续阅读 Mattila 的
《Geometry of Sets and Measures in Euclidean Spaces》\cite{Mattila1995Geometry}；
这里已把所需的紧集概率测度序列紧性和开集不等式在本节内证明；
\chapref{b1:ch:12}则把这些结论放到更一般的测度弱收敛框架中。

\subsection{Hausdorff 维数下界}
\label{b1:subsec:150403}

\begin{corollary}\label{b1:cor:frostman-dimension}
对非空紧集 \(K\subset\mathbb R^d\)，
\[
 \dim_{\mathrm H}K=
 \sup\{\,s\geq0\mid K\text{ 上存在增长指数为 }s\text{ 的 Frostman 测度}\,\}.
\]
\end{corollary}
\begin{proof}
若存在增长指数为 \(s\) 的测度，质量分布原理给出 \(\dim_{\mathrm H}K\geq s\)。
反之，若 \(0<s<\dim_{\mathrm H}K\)，则 \(\mathcal H^s(K)=\infty\)，
由 Frostman 引理可构造所需测度。
当 \(\dim_{\mathrm H}K=0\) 时，Dirac 测度给出指数 \(0\)，也得到等式。
\end{proof}

这不是说临界指数一定可以取到。设 \(s_0=\dim_{\mathrm H}K>0\)，
在 \(s=s_0\) 时，能否找到相应测度，恰好取决于
\(\mathcal H^{s_0}(K)\) 是否为正；维数本身并不决定这个数值。
在三分 Cantor 集上，前面构造的等分质量恰好实现临界指数。
对于别的集合，先证明所有 \(s<s_0\) 都有增长估计，已经足以得到维数下界，
不应为获得一个并非必要的端点估计而停下整个论证。

作为应用，设 \(F:K\rightarrow\mathbb R^m\) 是双 Lipschitz 映射，
\(\mu\) 是 \(K\) 上的一个 \(s\) 指数 Frostman 测度。
若逆映射的 Lipschitz 常数为 \(L\)，并且 \(B(y,r)\cap F(K)\ne\varnothing\)，
任选其中一点 \(F(x_0)\)，便有
\[
 F^{-1}\bigl(B(y,r)\cap F(K)\bigr)\subset B(x_0,2Lr).
\]
所以推前测度 \(F_\#\mu\) 仍满足 \(s\) 次增长估计。
这从测度构造一侧再次解释了双 Lipschitz 不变性：
距离只改变固定倍数时，质量增长的指数不会改变。

\subsection{能量方法初步}
\label{b1:subsec:150404}

逐球估计要求控制所有中心和尺度，有时不便直接验证。
另一种办法是把点对之间的距离作一次非负积分。

\begin{definition}\label{b1:def:riesz-energy}
设 \(t>0\)，\(\mu\) 为 \(\mathbb R^d\) 上的有限正 Borel 测度。
定义其 \term[riesz-nengliang]{Riesz 能量}[Riesz energy] 为
\[
 I_t(\mu)\coloneqq\iint_{\mathbb R^d\times\mathbb R^d}|x-y|^{-t}\,d\mu(y)\,d\mu(x),
\]
并约定对角线上的核值为 \(+\infty\)。
相应的非负函数
\[
 U_t^\mu(x)\coloneqq\int_{\mathbb R^d}|x-y|^{-t}\,d\mu(y)
\]
称为这里的 \(t\) 次势。能量与势都允许取无穷值。
\end{definition}
\symidx{\(I_t(\mu)\)：Riesz 能量}
\symidx{\(U_t^\mu\)：距离负幂核的势}

核是非负 Borel 函数，故 Tonelli 定理保证上述积分有意义，
也保证 \(U_t^\mu\) 可测。这里用距离负幂 \(t\) 作为下标；
其他书中若按势算子的阶数编号，可能使用核 \(|x-y|^{\alpha-d}\)，
其参数应对应为 \(t=d-\alpha\)，不能只按字母名称比较。

\begin{proposition}\label{b1:prop:frostman-finite-energy}
若非零有限正测度 \(\mu\) 满足
\(\mu\bigl(B(x,r)\bigr)\leq A r^s\) 对所有 \(x,r\) 成立，
则对 \(0<t<s\) 有 \(I_t(\mu)<\infty\)。
\end{proposition}
\begin{proof}
记 \(M=\mu(\mathbb R^d)\)。由
\[
 u^{-t}=\int_u^\infty t r^{-t-1}\,dr
\]
及 Tonelli 定理，包括 \(u=0\) 的无穷值情形，得到
\begin{equation}\label{b1:eq:energy-ball-integral}
 I_t(\mu)=t\int_0^\infty r^{-t-1}
       \int_{\mathbb R^d}\mu\bigl(B(x,r)\bigr)\,d\mu(x)\,dr.
\end{equation}
在 \(0<r<1\) 用增长估计，在 \(r\geq1\) 用总质量估计，即得
\[
 I_t(\mu)\leq tAM\int_0^1r^{s-t-1}\,dr+
 tM^2\int_1^\infty r^{-t-1}\,dr
 =\frac{tAM}{s-t}+M^2<\infty.
\]
\end{proof}

\begin{proposition}\label{b1:prop:finite-energy-dimension}
设非零有限正测度 \(\mu\) 集中在 Borel 集 \(E\) 上。
若 \(I_t(\mu)<\infty\)，则 \(\dim_{\mathrm H}E\geq t\)。
\end{proposition}
\begin{proof}
由能量有限，\(U_t^\mu(x)<\infty\) 对 \(\mu\)-几乎处处的 \(x\) 成立。
可选正整数 \(n\)，使 Borel 集
\[
 E_n=\{\,x\in E\mid U_t^\mu(x)\leq n\,\}
\]
有正质量，令 \(\nu=\mu|_{E_n}\)。
若 \(B(z,r)\cap E_n\ne\varnothing\)，取 \(x\) 属于此交集。
对 \(y\in B(z,r)\) 有 \(|x-y|<2r\)，故
\[
 \nu\bigl(B(z,r)\bigr)
 \leq(2r)^t\int_{B(z,r)}|x-y|^{-t}\,d\mu(y)
 \leq 2^t n r^t.
\]
若交集为空，左边为零，同一估计仍成立。
对正质量测度 \(\nu\) 使用质量分布原理，便得
\(\mathcal H^t(E_n)>0\)，因而 \(\dim_{\mathrm H}E\geq t\)。
\end{proof}

有限能量先给出势几乎处处有限，再限制到势有统一上界的部分；
这个限制步骤不能省略为“所以原测度具有统一增长估计”。
而且有限能量必定排除原子：若某点质量为 \(b>0\)，
对角线上的该点对就有质量 \(b^2\)，使能量为无穷。

结合以上两命题与 Frostman 引理，对非空紧集 \(K\) 有
\[
 \dim_{\mathrm H}K=
 \sup\Bigl(\{0\}\cup
 \{\,t>0\mid \text{存在集中在 }K\text{ 上的概率测度 }\mu,\ I_t(\mu)<\infty\,\}\Bigr).
\]
其中下界使用有限能量命题；上界的反向逼近则对
\(0<t<s<\dim_{\mathrm H}K\) 先构造 \(s\) 指数 Frostman 测度。
公式中的严格不等式 \(t<s\) 有实质作用，在端点把积分
\(\int_0^1r^{-1}\,dr\) 当成有限会得出错误结论。
Falconer 的《Fractal Geometry: Mathematical Foundations and Applications》%
\cite{Falconer2014Fractal}可作为继续学习维数估计与分形应用的伴读；
本节只建立了距离负幂核所需的基本积分估计，尚未引入一般容量理论。
